\subsection{特征 0 时滤过双代数的结构}

本号中我们继续假设 K 是特征 0 的域。

若 E 是双代数，典范单射 \(P(E) \rightarrow E\) 可以扩张为双代数态射 \(f_{E}: U(P(E)) \rightarrow E\)（第 4 号，命题 3）。

定理 1. 设 E 为余交换双代数。

\begin{enumerate}
\def\labelenumi{(\alph{enumi})}
\item
  双代数态射 \(f_{\mathbf{E}} \colon \mathrm{U}(\mathrm{P}(\mathbf{E})) \to \mathbf{E}\) 是单射。
\item
  若 E 上存在与其双代数结构相容的滤过（第 3 号，定义 2），则态射 \(f_{\mathrm{E}}\) 是同构。
\end{enumerate}

（在情形 (b)，双代数 E 因此等同于其本原元素李代数的包络双代数。）

设 \(c_{\mathbb{E}}\)（相应地，\(\varepsilon_{\mathbb{E}}\)）为 E 的余积（相应地，余单位）。我们记 \(\mathfrak{g} = \mathrm{P}(\mathrm{E})\)；设 \((e_i)_{i \in \mathrm{I}}\) 为向量 K-空间 \(\mathfrak{g}\) 的基，其中指标集 I 是全序的，并设 \((e_\alpha)_{\alpha \in \mathbf{N}^{(I)}}\) 为前一号中引入的基。我们记 \(\mathbf{X}_\alpha = f_{\mathbb{E}}(e_\alpha)\)（对 \(\alpha \in \mathbf{N}^{(I)}\)）。由 (15) 与 (16)，我们有：

\begin{enumerate}
\def\labelenumi{(\arabic{enumi})}
\setcounter{enumi}{16}
\tightlist
\item
\end{enumerate}

\[
\varepsilon_ {E} (X _ {0}) = 1, \quad \varepsilon_ {E} (X _ {\alpha}) = 0 \quad \text { for } | \alpha | \geqslant 1,\tag{18}
\]

\[
c _ {E} (X _ {\alpha}) = \sum_ {\beta + \gamma = \alpha} X _ {\beta} \otimes X _ {\gamma} \quad \text { for } \alpha \in \mathbf {N} ^ {(I)},
\]

因为 \(f_{E}\) 是余代数态射。

我们证明 \(f_{E}\) 是单射。这由下述引理得出：

引理 2. 设 V 为向量空间，E 为余代数，\(f\colon\mathsf{S}(\mathrm{V})\to\mathrm{E}\) 为余代数态射。若 \(f\) 在 \(\mathsf{S}^{\circ}(\mathrm{V})+\mathsf{S}^{1}(\mathrm{V})\) 上的限制是单射，则 \(f\) 是单射。

设 \(n \geqslant 0\)；我们记 \(\mathbb{S}_n = \sum_{i \geqslant n} \mathbb{S}^i(\mathbf{V})\)，用 \(c_s\) 表示 \(\mathbb{S}(\mathbf{V})\) 的余积，并对 \(n\) 用归纳法证明 \(f \mid \mathbb{S}_n\) 是单射。由于断言对 \(n = 0\) 与 \(n = 1\) 是平凡的，我们设 \(n \geqslant 2\)，并设 \(u \in \mathbb{S}_n\) 满足 \(f(u) = 0\)。那么

\[
\begin{array}{r l} {0 = c _ {\mathbf {E}} (f (u))} & {= (f \otimes f) (c _ {\mathbf {s}} (u))} \\ & {= f (u) \otimes 1 + 1 \otimes f (u) + (f \otimes f) (c _ {\mathbf {s}} ^ {+} (u))} \\ & {= (f \otimes f) (c _ {\mathbf {s}} ^ {+} (u)).} \end{array}
\]

由于 \(c_{\mathbb{S}}^{+}(u) \in \mathbb{S}_{n-1} \otimes \mathbb{S}_{n-1}\)，由 (11) 与归纳假设可知 \(u\) 是 \(\mathbb{S}(\mathrm{V})\) 的本原元素，因而是次数 1 的（第 5 号，注 4），从而由于 \(f \mid \mathbb{S}^1(\mathrm{V})\) 是单射而为零。

特别地由此得出族 \((X_{\alpha})\) 是自由的。

我们证明：\(f_{\mathbb{E}}\) 当 \(\mathrm{E}\) 带有与其双代数结构相容的滤过时是满射。设 \((\mathrm{E}_n)_{n\geqslant 0}\) 为这样一个滤过，记 \(\mathrm{E}_n^+ = \mathrm{E}_n\cap \mathrm{Ker}(\epsilon_{\mathbb{E}})\)。我们对 \(n\) 用归纳法证明 \(\mathrm{E}_n^+\) 包含在 \(f_{\mathbb{E}}\) 的像中。由于 \(\mathrm{E} = \mathrm{K}.1 + \bigcup_{n > 0}\mathrm{E}_n^+\)，这将蕴含 \(f_{\mathbb{E}}\) 的满射性。断言对 \(n = 0\) 是平凡的，对 \(n = 1\) 由第 3 号命题 6 的推论得出；以下设 \(n\geqslant 2\)，并设 \(x\in \mathrm{E}_n^+\)。由第 3 号命题 6，

\[
c _ {\mathrm{E}} ^ {+} (x) \in \sum_ {i = 1} ^ {n - 1} \mathrm{E} _ {i} ^ {+} \otimes \mathrm{E} _ {n - i} ^ {+}
\]

且由归纳假设，存在标量 \(\lambda_{\alpha, \beta}\)，其中 \(\alpha, \beta\) 属于 \(\mathbf{N}^{(I)}\)，除有限多个外均为零，使得

\[
c _ {\mathbb {E}} ^ {+} (x) = \sum_ {\alpha , \beta \neq 0} \lambda_ {\alpha , \beta} \mathrm{X} _ {\alpha} \otimes \mathrm{X} _ {\beta}.\tag{19}
\]

于是由公式 (18)

\[
(c _ {\mathrm{E}} ^ {+} \otimes \operatorname{Id} _ {\mathrm{E}}) (c _ {\mathrm{E}} ^ {+} (x)) = \sum_ {\alpha , \beta , \gamma \neq 0} \lambda_ {\alpha + \beta , \gamma} \mathbf {X} _ {\alpha} \otimes \mathbf {X} _ {\beta} \otimes \mathbf {X} _ {\gamma}
\]

\[
(\mathrm{Id} _ {\mathtt {E}} \otimes c _ {\mathtt {E}} ^ {+}) (c _ {\mathtt {E}} ^ {+} (x)) = \sum_ {\alpha , \beta , \gamma \neq 0} \lambda_ {\alpha , \beta + \gamma} \mathbf {X} _ {\alpha} \otimes \mathbf {X} _ {\beta} \otimes \mathbf {X} _ {\gamma}.
\]

由第 1 号命题 3 以及诸 \(X_{\alpha}\) 的线性无关性，得出

\[
\lambda_ {\alpha + \beta , \gamma} = \lambda_ {\alpha , \beta + \gamma} \quad \text { for } \alpha , \beta , \gamma \text { in } \mathbf {N} ^ {(I)} - \{0 \}.\tag{20}
\]

另一方面，余积 \(c_{E}\) 是余交换的；与上面同样的论证蕴含

\[
\lambda_ {\alpha , \beta} = \lambda_ {\beta , \alpha} \quad \text { for } \alpha , \beta \text { in } \mathbf {N} ^ {(I)} - \{0 \}.\tag{21}
\]

假设存在标量族 \((\mu_{\alpha})\)（\(|\alpha| \geqslant 2\)），使得

\[
\mu_ {\alpha + \beta} = \lambda_ {\alpha , \beta} \quad \text { for } \alpha , \beta \text { in } \mathbf {N} ^ {(I)} - \{0 \}.\tag{22}
\]

那么

\[
c _ {\mathrm{E}} ^ {+} (x) = \sum_ {\alpha , \beta \neq 0} \mu_ {\alpha + \beta} \mathrm{X} _ {\alpha} \otimes \mathrm{X} _ {\beta} = \sum_ {| \gamma | \geqslant 2} \mu_ {\gamma} c _ {\mathrm{E}} ^ {+} (\mathrm{X} _ {\gamma})
\]

由公式 (18)，于是 \(y = x - \sum_{|\gamma| \geq 2} \mu_\gamma X_\gamma\) 是本原的，因而属于 \(P(E) \subset \operatorname{Im}(f_E)\)。从而

\[
x = y + \sum_ {| \gamma | \geqslant 2} \mu_ {\gamma} f _ {\mathrm{E}} (e _ {\gamma}) \in \operatorname{Im} (f _ {\mathrm{E}}).
\]

因此，只要证明了下述引理，证明即告完成：

引理 3. 若有限支集的标量族 \((\lambda_{\alpha, \beta})\)（其中 \(\alpha, \beta\) 属于 \(\mathbf{N}^{(I)} - \{0\}\)）满足关系 (20) 与 (21)，则存在有限支集的族 \((\mu_{\alpha})_{|\alpha| \geqslant 2}\)，使得 \(\mu_{\alpha + \beta} = \lambda_{\alpha, \beta}\) 对非零的 \(\alpha, \beta\) 成立。

只需证明

\[
\alpha + \beta = \gamma + \delta\tag{23}
\]

蕴含 \(\lambda_{\alpha, \beta} = \lambda_{\gamma, \delta}\) 对非零的 \(\alpha, \beta, \gamma, \delta\) 成立。由 Riesz 分解引理（《代数》，第 VI 章，第 1 节，第 10 号，定理 1），存在 \(\pi, \rho, \sigma, \tau\)（\(\mathbf{N}^{(I)}\) 中的），使得

\[
\alpha = \pi + \sigma , \quad \beta = \rho + \tau , \quad \gamma = \pi + \rho , \quad \delta = \sigma + \tau .
\]

设 \(\pi \neq 0\)；由于 \(\sigma + \beta = \rho + \delta\)，关系 (20) 蕴含

\[
\lambda_ {\alpha , \beta} = \lambda_ {\pi + \sigma , \beta} = \lambda_ {\pi , \sigma + \delta} = \lambda_ {\pi , \rho + \delta} = \lambda_ {\pi + \rho , \delta} = \lambda_ {\gamma , \delta}.
\]

反之，若 \(\pi = 0\)，则 \(\beta = \gamma + \tau\) 且 \(\delta = \alpha + \tau\)，由此

\[
\lambda_ {\alpha , \beta} = \lambda_ {\alpha , \gamma + \tau} = \lambda_ {\alpha + \tau , \gamma} = \lambda_ {\delta , \gamma}
\]

由 (20) 得出，而又由 (21) 有 \(\lambda_{\delta, \gamma} = \lambda_{\gamma, \delta}\)，故 \(\lambda_{\alpha, \beta} = \lambda_{\gamma, \delta}\)。

